\documentclass[11pt]{article}
\usepackage[margin=1in]{geometry}
\usepackage{amsmath,amssymb,amsthm,hyperref}
\newtheorem{theorem}{Theorem}\newtheorem{lemma}[theorem]{Lemma}
\newcommand{\E}{\mathbb E}\newcommand{\R}{\mathbb R}
\newcommand{\relu}{\operatorname{ReLU}}
\title{Bayesian feature selection among irrelevant inputs}
\author{ChatGPT-6 Astra\thanks{Independent AI research. Sole author. Verification scope and provenance are stated in the accompanying research collection.}}\date{18 November 2026; independent research draft}
\emergencystretch=3em
\hypersetup{pdftitle={Bayesian feature selection among irrelevant inputs},pdfauthor={ChatGPT-6 Astra}}
\begin{document}\maketitle

This extension keeps the target's intrinsic parity dimension fixed while
increasing the number of irrelevant Boolean coordinates. It shows that
the exceptional-neuron mechanism also describes a small learned feature,
rather than only a full-dimensional parity representation. Every cube
label is still observed; this is not a sample-complexity or OOD theorem.

\section{Fixed intrinsic parity, variable ambient dimension}
Let $r\ge2$ be even and $d\ge\max(r,6)$. On $X=\{-1,1\}^d$ let
$\chi(x)=\prod_{i=1}^r x_i$ and $M=2^d$. Use unit independent Gaussian
parameters and $F_N=N^{-1}\sum_j a_j\relu(b_j+w_j\cdot x)$.
Put
\[
 \rho_r=\frac{\binom{r-2}{(r-2)/2}}{2^{r-1}\sqrt r},\quad
 Z=\E_x\chi(x)a\relu(b+w\cdot x),\quad
 Q_{*,d}=e^{Z/\rho_r}P_{0,d}/M_{*,d},\quad m_{c,d}=\E_{Q_{*,d}}Z.
\]
The critical-tilt companion already proves that the endpoint law has
finite fourth parameter moment whenever the \emph{ambient} dimension
$d>5$. Its critical rays use only the $r$ relevant weights: bias and
irrelevant weights are zero.
Define
\[
 L_r=\begin{cases}4,&r=2,\\2^{r-1},&r\ge4.\end{cases}
\]
The notation $L_r$ here is a neuron count, unrelated to the ray-volume
constant in the orientation-law companion.

\begin{theorem}[Posterior specialization to the relevant coordinates]
Fix $d,r$ and $s>m_{c,d}$. Condition on $\chi(x)F_N(x)>s$ at every
cube input. As $N\to\infty$:
\begin{enumerate}
\item The ordinary empirical law tends to $Q_{*,d}$ with mean output
$m_{c,d}\chi$, while the whole output tends to $s\chi$.
\item Exactly $L_r$ neurons have norm above $\varepsilon\sqrt N$,
for every sufficiently small fixed $\varepsilon>0$. Each has squared
norm divided by $N$ tending to $2(s-m_{c,d})/(L_r\rho_r)$.
For $r\ge4$, one orientation per antipodal pair occurs, with the same
balancing equations as in the full-parity theorem. For $r=2$, all four
critical rays occur.
\item Their hidden biases and irrelevant incoming weights are
$O_{\Pr}(N^{-1/2})$. The relevant incoming vector differs from its
projection onto the selected equal-magnitude sign ray by
$O_{\Pr}(N^{-1/2})$. The difference between absolute readout and
incoming radial magnitude is $O_{\Pr}(1)$.
\item The learned feature kernel satisfies
$\mathcal K_N\chi\to s\rho_r\chi$.
The energy limits are $N^{-1}\sum a_j^2\to1+s/\rho_r$ and
$N^{-1}\sum\|(b_j,w_j)\|^2\to d+1+s/\rho_r$.
\item With $I_d(s)=s/\rho_r-\log M_{*,d}$, the evidence obeys
\[
 -\log P_{0,d}^N(H_N)=NI_d(s)+
 \left[\frac M2+\frac{L_r(d-2)}2\right]\log N+O_{d,r,s}(1).
\]
\end{enumerate}
The analogous above-critical square or logistic likelihood has the same
structure, and its logarithmic correction is $L_r(d-2)\log N/2$.
All constants and width thresholds may depend on fixed $d,r$.
\end{theorem}

\section{Why the earlier proof extends}
The coarse conditioning and truncated-moment arguments use only finite
fourth parameter moments, the full finite-domain output span, and finitely
many critical rays. All three hold here. The critical rays are
$(\sigma_z,v_z)$, with
$v_z=(0,z/\sqrt r,0_{d-r})$ and
$\sigma_z=\operatorname{sgn}(G(v_{\boldsymbol1}))\chi(z)$.
Their output vectors are constant on every irrelevant-coordinate fiber.

For $r\ge4$, the invertible even Fourier matrix on the relevant
antipodal quotient fixes the positive mass in every pair exactly as
before. The canonical hemisphere provides a representation using
$L_r=2^{r-1}$ independent feature columns and equal amplitudes.
For $r=2$, a critical feature is nonzero on exactly one of the four
relevant input patterns, with magnitude $\sqrt2$ and the correct
parity sign. Thus all four rays are necessary, their four feature
columns are independent, and each amplitude is
$(s-m_{c,d})/(4\rho_r)$. This is why merely substituting
$2^{r-1}$ into the two-bit case would be wrong.

Near a critical ray, the transverse incoming space has dimension $d$.
The coordinates $c=AR$, $v=A-R$, $t=\sqrt c\,u$ therefore have density
bounded by
$Cc^{-(d+1)/2}e^{-c_0(v^2+\|t\|)}$.
For $L$ selected macroscopic neurons whose feature columns have rank $R$,
the exact-count companion's coarea and Gaussian-component estimates give
weighted union mass at most
\[
 C N^{-\{L(d-3)+M+R\}/2+\nu},\qquad
 \nu=0\ \text{for hard margins},\quad \nu=M/2\ \text{for soft likelihoods}.
\]
The coarse geometry requires $L\ge L_r$ and $R\ge L_r$.
The matching lower construction selects $L_r$ independent rays, allows
$O(\sqrt N)$ amplitude fluctuations, and counts the $(N)_{L_r}$ index
assignments. It has exactly the asserted power. Every extra macroscopic
neuron loses at least $N^{-(d-3)/2}$, proving the exact count.
The same invertible core equations give the equal limiting masses.

The exponential envelope in $(v,t)$, with the same normalized upper
bound, makes those variables tight under the posterior. Since
$c/N\to(s-m_{c,d})/(L_r\rho_r)>0$, this proves
$u=t/\sqrt c=O_{\Pr}(N^{-1/2})$ and $A-R=O_{\Pr}(1)$.
The learned-kernel and energy arguments use
$G(U)\relu(U\cdot\bar x)=\rho_r\Phi_x$ and
$a^2=\|U\|^2=Z/\rho_r$ on each critical ray, and the same quadratic
truncation argument for the ordinary population. These prove the
remaining assertions. No new assumption about irrelevant coordinates
is hidden in the conditioning step: all $2^d$ output constraints enter
the full-dimensional Gaussian-component calculation.

\section{The two-bit result also holds in ambient dimensions four and five}
For $r=2$, every conclusion of the first theorem extends to $d=4,5$
whenever $s>m_{c,d}$, despite the infinite full output covariance at
those critical laws. The same extension holds above the soft-likelihood
transition. The distinction is that the four critical feature columns
span the \emph{entire} space of critical-ray outputs. Their amplitudes
can absorb the non-Gaussian bulk fluctuations in that space.

Here are the modifications needed in the proof. Let $E$ be that
four-dimensional feature span and let $P_\perp$ project onto $E^\perp$.
Near a critical ray the output is $cf_z+e$, with
$\|e\|\le C\|t\|$ and density bounded by
$Cc^{-(d+1)/2}e^{-c_0(v^2+\|t\|)}$. Thus $P_\perp\Phi$ has all
moments, and indeed a small exponential moment. Outside the ray cones
the critical density has a Gaussian quadratic gap. The projected output
has positive definite covariance on $E^\perp$ and a convolution block
with an isotropic Gaussian component there. Its sum therefore has the
same centered local lower and upper density bounds
$N^{-(M-4)/2}$ as in the earlier Gaussian-component lemma.
The full output still has a finite $q$th moment for some $q>1$:
choose $1<q<(d-1)/2$.

For a direct polynomial evidence lower bound, reserve one neuron on each
of the four rays. Let the other neurons' projected output fall in a fixed
small box about zero. Their parallel output differs from its mean by
$o(N)$ while retaining a fixed fraction of this projected-box probability.
To justify the latter assertion, use the Gaussian-component decomposition:
the residual projected sum is $O(\sqrt N)$ on an event of fixed positive
probability, by Chebyshev; its parallel sum is $o(N)$ in probability by
the finite first moment and the law of large numbers. These events
intersect with fixed positive probability. Conditional on the residuals,
the Gaussian part supplies the projected box with probability
$\Omega(N^{-(M-4)/2})$, and its parallel fluctuation is $O(\sqrt N)$.

Given such a bulk realization, integrate the four reserved amplitudes
near the unique values that cancel its parallel fluctuation and produce
signed total outputs between $Ns+1$ and $Ns+3$. Each amplitude is
$N(s-m_{c,d})/(4\rho_r)+o(N)>0$. A fixed-size four-dimensional amplitude
box has density at least $cN^{-4(d+1)/2}$; restrict the other local
coordinates to small fixed boxes. Counting the $N^4$ assignments, with
the bounded-overlap argument below, gives
\[
 Z_N^{\rm hard}\ge cN^{-[M/2+2d-4]}.
\]
For the soft likelihood allow projected and final output fluctuations
of order $\sqrt N$, and integrate an amplitude box of side $\sqrt N$.
This gives $Z_N^{\rm soft}\ge cN^{-(2d-4)}$.

The lower-bound overlap argument and the posterior coarse geometry need
only a parameter moment of some order $p>2$. Such a moment exists since
$p<d-1$ is allowed. Markov and the index-set union bound give
\[
 Q_*^N\{L_N>D\}\le C_DN^{-(p/2-1)(D+1)}.
\]
Choose the fixed count cutoff $D$ large enough to beat the evidence
power plus the four index factors. On $L_N\le D$, the assignment
events overlap at most $(D)_4$ times. This proves the stated lower bound
without assuming disjoint assignments. It also bounds the posterior
macroscopic count by a fixed constant after polynomial domination.

For the collective contribution of small neurons, replace the second-
moment Taylor estimate by
\[
 |e^x-1-x|\le C_q|x|^q e^{|x|},\qquad 1<q\le2.
\]
With truncation $\|\theta\|\le\varepsilon\sqrt N$ and
$\lambda=A\log N/N$, the log moment of the centered sum is at most
$C N\lambda^q e^{C\lambda\varepsilon^2N}=o(1)$ after taking
$\varepsilon$ sufficiently small. This gives any prescribed polynomial
concentration power as before. The same argument applies to quadratic
parameter and kernel observables, using $2q<d-1$.
Angular concentration and the original-prior energy bound are unchanged.

All four critical rays must then be present with their positive masses.
For $L\ge4$ selected large neurons, the feature-column rank is exactly
four. The coarea upper bound, using the good Gaussian component of the
remaining output sum, is the same as before, with
$N^{-\{L(d-3)+M+4\}/2+\nu}$. It uses no moment bound for the residual
translation. Every extra neuron loses $N^{-(d-3)/2}$, a positive power
for both $d=4$ and $d=5$. This proves exact count four and matching
evidence powers. The remaining specialization and kernel conclusions
follow from the $q$-moment truncation argument. No full-covariance central
limit theorem has been invoked where it is false.

This extension does not settle full four-bit parity in ambient dimension
four. There the critical-ray span is larger than a minimal orientation
set, and non-Gaussian fluctuations remain transverse to its amplitudes.
The Gaussian covariance determinant formula in the orientation note
must not be applied in that case.

\section{The ordinary law approaches the untrained prior}
\begin{theorem}[Irrelevant dimensions suppress the critical ordinary signal]
For every fixed even $r\ge2$, as $d\to\infty$,
\[
 M_{*,d}-1=O_r(d^{-1/2}),\qquad m_{c,d}=O_r(d^{-1/2}),\qquad
 \operatorname{TV}(Q_{*,d},P_{0,d})=O_r(d^{-1/4}).
\]
Thus every fixed $s>0$ is above the critical mean for all sufficiently
large ambient dimensions. These rates are convenient upper bounds,
not claims of sharp dimension exponents.
\end{theorem}
\begin{proof}
Fix $d_0=\max(r,4)$ relevant-plus-base coordinates. Let
$G_{d_0}(b,w)$ denote their parity feature and set
\[
 \psi_0(g)=e^{g^2/(2\rho_r^2)}-1,\qquad
 \psi_1(g)=g^2e^{g^2/(2\rho_r^2)}.
\]
Both functions are convex. We first note that
\[
 J_i=\int_{\R}\E_{w\sim N(0,I_{d_0})}
                  \psi_i(G_{d_0}(b,w))\,db<\infty,
 \qquad i=0,1.\tag{1}
\]
Here the bias is integrated against Lebesgue measure, not a Gaussian.
The integrand is zero if $|b|>\sum_j|w_j|$, because the ReLU becomes
affine or zero on the entire cube. The sharp minimum-weight inequality
gives $|G|\le\rho_r\sqrt r\min_{j\le r}|w_j|\le\rho_r\|w\|$.
Equality forces equal relevant magnitudes and zero irrelevant weights;
for these weights the one-dimensional parity hinge sum has its unique
absolute maximum at bias zero. The latter follows from its successive
binomial coefficients: for even $r$ the central coefficient is strictly
larger than its neighbors. Consequently the Gaussian-weight exponent
$\|w\|^2-G^2/\rho_r^2$ has a strictly positive quadratic gap away
from the same finitely many critical rays, on the compact normalized
region $|b|\le\sqrt{d_0}\|w\|$.

Near a ray, deleting the bias-square term from the earlier corner
calculation still leaves a quantity comparable to $R\|u\|$ after
shrinking the fixed cone. The first-moment integral is bounded by
$C\int_1^\infty R^{2-d_0}dR$, which is finite since $d_0>3$;
the zeroth-moment integral is easier. Away from the cones the quadratic
gap gives integrability, and bounded parameter regions cause no issue.
This proves (1).

Average over the remaining Boolean inputs before applying $\psi_i$.
Jensen's inequality gives
\[
 \E_{U\sim N(0,I_{d+1})}\psi_i(G_d(U))
 \le \E_{B\sim N(0,d-d_0+1),\,w\sim N(0,I_{d_0})}
                       \psi_i(G_{d_0}(B,w))
 \le\frac{J_i}{\sqrt{2\pi(d-d_0+1)}}.
\]
The first line uses only that, for every fixed irrelevant sign pattern,
the shifted bias has the displayed Gaussian variance and is independent
of the base weights. The bound for $i=0$ is the assertion about $M_*-1$.
Since $\rho_rM_{*,d}m_{c,d}=\E\psi_1(G_d)$ and $M_{*,d}\ge1$,
it also proves the mean bound. Finally,
$\operatorname{KL}(Q_{*,d}\|P_{0,d})=m_{c,d}/\rho_r-\log M_{*,d}
\le m_{c,d}/\rho_r$, and Pinsker gives the total-variation estimate.
\end{proof}

\section{The initial kernel loses the same fixed target feature}
For fixed even $r$, the unit-Gaussian initial ReLU kernel has target
parity eigenvalue
\[
 \lambda_r(d)\sim C_r d^{1-r},\qquad
 C_r=\frac{(r-2)!\binom{r-2}{(r-2)/2}}{2\pi\,2^{r-2}}>0.
\]
To see this, write its covariance as
$(d+1)g((1+x\cdot y)/(d+1))$, where
\[
 g(c)=\frac{\sqrt{1-c^2}+(\pi-\arccos c)c}{2\pi}.
\]
Taking $r$ finite differences yields exactly
\[
 \lambda_r(d)=(d+1)^{1-r}\E\,g^{(r)}
 \left(\frac{1+\sum_{i=r+1}^d\varepsilon_i+
                         \sum_{i=1}^r T_i}{d+1}\right),
\]
where the $\varepsilon_i$ are independent uniform signs and
$T_i$ independent uniform variables on $[-1,1]$.
On $|\sum\varepsilon_i|\le d/2$, the derivative is bounded and its
argument tends to zero in probability, so the expectation tends to
$g^{(r)}(0)$. The complementary sign event has exponentially small
probability. Its worst conditional $T$ integral is only
$O_r(d^{r-3/2})$: the endpoint singularity is
$(1-|c|)^{3/2-r}$, and integration over the $r$ continuous coordinates
makes it integrable. Thus that part vanishes. Finally expanding
$g''(c)=(2\pi\sqrt{1-c^2})^{-1}$ gives the displayed constant.

For the especially simple target $x_1x_2$, the initial eigenvalue is
asymptotic to $1/(2\pi d)$, while the learned value is
$s/(2\sqrt2)$. Exactly four exceptional neurons carry the missing
signal. Their irrelevant weights tend to zero as $N^{-1/2}$, even as
the typical neuron law approaches its original Gaussian prior in the
subsequent large-$d$ limit. This is explicit specialization to the
relevant coordinates. The target's four-neuron representation is fixed
throughout; the experiment still observes all cube labels and gives no
claim about discovering those coordinates from a small sample.
\section{Numerical check of the dimension crossover}
Independent importance integration for the fixed two-bit target gives:
\[
\begin{array}{c|rrrrrr}
 d&4&5&6&8&10&12\\\hline
 \widehat m_{c,d}&1.5471&0.6682&0.3870&0.1807&0.1072&0.0735
\end{array}
\]
The first two entries use 50 million proposals each; the next three use
10 million each, and the last uses five million. Their empirical Monte
Carlo standard errors are respectively about
$0.00090,0.00044,0.00068,0.00046,0.00042,0.00048$.
These are exploratory floating-point estimates, not deterministic
integration certificates. They suggest that margin $s=1/2$ crosses the
critical value between $d=5$ and $d=6$. The analytic result is the eventual
transition as $d$ increases, not a certified sharp transition dimension.

For this two-bit sampler the finite cube average has a stable closed
form. If the relevant weight magnitudes are $t,t+v$, $v\ge0$, and
$B=b+\sum_{i>2}w_ix_i$, then
\[
 G=\frac t2-\frac14\E_{x_{>2}}
                 \min\{\max(|B|-v,0),2t\}.
\]
This follows by pairing the two positive and two negative parity inputs
and using $\relu(B+A)+\relu(B-A)=B+\max(|B|,A)$.
It avoids cancellation at large critical weights and reduces the
numerical cube sum by a factor of four. The sampler also checks the exact
incoming-minus-readout energy identity $d$.
\section{Changing the Gaussian layer scales}
The preceding results are not tied to unit incoming variance. Let
$a=\tau_a\widetilde a$, $U=\tau_u\widetilde U$, with independent
standard Gaussian standardized parameters and fixed $\tau_a,\tau_u>0$.
Positive homogeneity gives $\Phi=\tau_a\tau_u\widetilde\Phi$ exactly.
Consequently the physical-margin posterior at $s$ is precisely the
standardized posterior at $\widetilde s=s/(\tau_a\tau_u)$. Its critical
ordinary margin is $\tau_a\tau_u m_{c,d}$ and its critical dual parameter
is $(\tau_a\tau_u\rho_r)^{-1}$. All exact-count, orientation, and
localization assertions transfer by this invertible linear map.

Since the incoming-feature kernel scales by $\tau_u^2$, its learned
target eigenvalue is
\[
            \lambda_\chi(\mathcal K)\longrightarrow
                        s\rho_r\frac{\tau_u}{\tau_a}.
\]
In the above-critical regime write
$\delta=s-\tau_a\tau_u m_{c,d}>0$. Each of the $L_r$ exceptional
neurons has
\[
 \frac{a_j^2}{N}\longrightarrow
       \frac{\delta\tau_a}{L_r\rho_r\tau_u},\qquad
 \frac{\|U_j\|^2}{N}\longrightarrow
       \frac{\delta\tau_u}{L_r\rho_r\tau_a}.
\]
These formulas concern fixed dimension and fixed layer scales before
the width limit.

For the usual incoming scaling $\tau_u=d^{-1/2}$ and $\tau_a=1$,
two-bit parity therefore has prior target-kernel eigenvalue asymptotic
to $1/(2\pi d^2)$ and learned eigenvalue $s/(2\sqrt{2d})$.
Their ratio is asymptotic to $(\pi s/\sqrt2)d^{3/2}$.
The four exceptional incoming vectors still discard the irrelevant
coordinates, and the physical critical margin tends to zero at least
as fast as $d^{-1}$. This is an iterated dimension--width corollary;
it does not give a sufficient finite width uniform in $d$.
\end{document}
